\chapter{Simulation Study}
\section{Design and comparisons}
The completed study contains six data-generating processes and
\SimReplications{} paired replications per scenario, giving 600 replications
in total. Each replication fits Cox regression and a log-rank random survival
forest twice to the same retained sample: once with entry times zeroed and
once with observed entry times supplied. Both phases use left-truncated
data; zeroing entry does not restore subjects excluded by truncation.

Each replication also independently fits the ordinary versions of Cox and
RSF to verify agreement with their zero-entry counterparts. Thus, six fits
are performed per replication. The comparison of interest remains the
within-family change from entry-blind to entry-aware fitting.

\begin{table}[htbp]
\centering\small
\setlength{\tabcolsep}{4pt}
\begin{tabular}{cll}
\toprule
Scenario & Event-time family & Linear Cox specification\\
\midrule
1 & Weibull & Correctly specified PH\\
2 & Covariate-dependent Gompertz shape & Non-proportional hazards\\
3 & Exponential & Correctly specified PH\\
4 & Exponential with nonlinear mean & Misspecified link; PH holds\\
5 & Covariate-dependent Gamma shape & Non-proportional hazards\\
6 & Lognormal & Non-proportional hazards\\
\bottomrule
\end{tabular}
\caption{Overview of the six scenarios. PH denotes proportional hazards.
Scenario 4 illustrates a nonlinear log-hazard rather than a failure of
proportional hazards. Complete DGP formulas and source-paper comparisons
remain part of the final methods write-up.}
\label{tab:dgp}
\end{table}

Entry is drawn independently of the initial event-time and covariate
draws. A subject is retained only if it survives to entry. Censoring is
generated as a residual duration after entry, using scenario-specific
distributions and administrative cutoffs. This sampling sequence differs
from drawing an absolute censoring time before selection. Source-population
independence does not imply independence within the selected cohort.

\section{Run size and model settings}
Each replication uses \SimTrainN{} retained training subjects from a
candidate pool of 50,000, \SimTrees{} trees per forest, and
\SimTestN{} independently drawn population test covariates. Survival curves
are evaluated on the preserved 100-point scenario-specific grid.
Table~\ref{tab:hyperparameters} gives the forest settings. The existing
setup files and master seeds were preserved; no retuning was performed.
The two phases use matching seeds and hyperparameters.

% Generated from the preserved run configuration.
\begin{table}[htbp]
\centering
\begin{tabular}{crr}
\toprule
Scenario & mtry & nodesize \\
\midrule
1 & 3 & 16 \\
2 & 3 & 16 \\
3 & 5 & 6 \\
4 & 3 & 6 \\
5 & 5 & 6 \\
6 & 5 & 6 \\
\bottomrule
\end{tabular}
\caption{Preserved forest hyperparameters. Every scenario uses \SimTrees{} trees and \SimTrainN{} retained training subjects per replication. No retuning was performed for this refresh.}
\label{tab:hyperparameters}
\end{table}


The refreshed run uses package version \texttt{3.9.7} with the LT changes
at commit \RunCommit{}. It includes the upstream updates merged into the
project and the subsequent bug fixes. Installed R and native-library
fingerprints are recorded with every result, and workers verify that they
use the matching installation. Completed replications are checkpointed
individually. The six 100-replication result objects were completed on
7 November 2026.

\section{Evaluation samples and curve accuracy}
For $m=\SimTestN$ population test covariates and $G=100$ grid times,
\[
 \operatorname{MAD}=\frac{1}{Gm}\sum_{g=1}^{G}\sum_{i=1}^{m}
 \left|\widehat S(t_g\mid X_i)-S(t_g\mid X_i)\right|.
\]
This is an equally weighted grid average, not an interval-length-weighted
time integral. The test sample supplies covariates and analytic survival
truth; it does not supply simulated validation outcomes.

Concordance, Brier/CRPS, and AUC use observed training outcomes and real
entry times in both fitting phases. RSF contributes OOB curves, while Cox
contributes in-sample curves. Within a model family this holds the scoring
target fixed while varying the entry information used in fitting.
Across families the evaluation procedures differ, limiting direct
interpretation of outcome-score rankings. The pooled censoring model and
its limitations are described in Chapter 3.

\section{Paired differences and Monte Carlo uncertainty}
For replication-level errors $E_{0r}$ and $E_{Lr}$, define
\[
 D_r=E_{0r}-E_{Lr},\qquad
 \overline D=\frac{1}{R}\sum_{r=1}^R D_r,\qquad
 \operatorname{MCSE}(\overline D)=\frac{s_D}{\sqrt R}.
\]
All four reported metrics use error scales, so positive gains favor
entry-aware fitting. The tables report paired MCSEs. Descriptive
normal-approximation intervals can be formed as
$\overline D\pm1.96\operatorname{MCSE}(\overline D)$; they quantify
Monte Carlo uncertainty in the mean paired difference and are not
individual prediction intervals.

Percentage reduction is
$100(\overline E_0-\overline E_L)/\overline E_0$, a ratio of mean errors.
It is not the mean of replication-specific percentage changes.
For concordance and AUC, this is a relative change in $1-C$ or
$1-\overline{\mathrm{AUC}}$, not a percentage-point increase in C or AUC.
